\section{Computational and High Dimensional Considerations}
\label{sec:computational_considerations}

As illustrated in Figure \ref{fig:sbc_workflow}, SBC can be computationally expensive, as it requires repeated posterior inference for many synthetic datasets.
This cost can be mitigated through parallelisation, since each inference run is independent.
Despite the high computational cost, performing SBC early in the modelling process allows for detection of coding mistakes or conceptual errors before substantial resources are spent on large-scale analyses.
For this reason, it is recommended that SBC be considered as part of the Bayesian model development process, rather than an afterthought, as emphasised by \textcite{gelman_2020}.

These computational challenges become even more pronounced in high-dimensional models, where a large number of parameters must be assessed simultaneously.
Applying SBC independently to each parameter is conceptually straightforward but has a huge cost.
To reduce this cost, the workflow in Figure \ref{fig:sbc_workflow} can be extended to high dimensional models by just taking $\theta$ as a $d$-dimensional vector.
For each component of $\theta$, rank statistics can be computed and stored across simulations, and their empirical distributions can then be checked for uniformity.
While this approach increases storage and computational demands, it provides a systematic way to diagnose inference quality even in complex, high-dimensional settings.
